\documentclass[11pt,a4paper]{article}
% Shared preamble for RuPhO English problem statements (match T1 2026 style).
% Use: \input{latex/rupho_en_preamble.tex} after \documentclass.
\usepackage[margin=1in]{geometry}
\usepackage{amsmath,amssymb}
\usepackage{graphicx}
\usepackage{float}
\usepackage{enumitem}
\usepackage{microtype}
\usepackage{caption}
\usepackage{tabularx}
\usepackage{hyperref}

\captionsetup{font=small,labelfont=bf,skip=6pt}

\setlist[itemize]{leftmargin=1.2cm}
\setlist[enumerate]{leftmargin=1.2cm}

% One outer box per subproblem: [id | statement | points] (no inner rules)
\newcommand{\problembox}[3]{%
  \par\addvspace{0.42em}%
  \noindent
  \setlength{\fboxsep}{7.5pt}%
  \setlength{\fboxrule}{0.95pt}%
  \fbox{%
    \begin{minipage}{\dimexpr\linewidth-2\fboxsep-2\fboxrule\relax}
    \begin{tabularx}{\linewidth}{@{}>{\raggedright\arraybackslash\bfseries}p{2.1em} @{\hspace{0.9em}} >{\raggedright\arraybackslash}X @{\hspace{0.9em}} >{\raggedleft\arraybackslash\bfseries}p{2.4em}@{}}
    #1 & #2 & #3
    \end{tabularx}
    \end{minipage}%
  }%
  \par\addvspace{0.32em}%
}

\title{\textbf{Plasmon resonance and SERS}}
\author{\normalsize Y18 (2018) --- English translation}
\date{}
\begin{document}
\maketitle

Surface-enhanced Raman spectroscopy (SERS) is one of the most striking optical phenomena discovered in the last 40 years. SERS is based on plasmon resonance—a strong increase in the electric field near metal nanobeads under certain conditions. To determine these conditions and understand how SERS works, it is necessary to learn how to describe the properties of a metal placed in an alternating electromagnetic field.

Let us recall that the properties of matter in an electric field obey the following relations: $$\vec{D} = \varepsilon_0 \varepsilon \vec{E}, \qquad \vec{D} = \varepsilon_0 \vec{E} + \vec{P},$$ where $\vec{E}$ and $\vec{D}$ are the electric field intensity and induction, respectively; $\varepsilon$ is the dielectric constant of the substance; $\vec{P}$ is the polarization vector (the dipole moment of a unit volume); and $\varepsilon_0 = 8.85\cdot 10^{-12}~\text{F}/\text{m}$ is the electric constant. At the interface between two media in the absence of free charges, the tangential component of the intensity vector and the normal component of the electric field induction vector are continuous. In the case of an alternating electric field, the dielectric constant of a substance (including metals) depends on the frequency of the electromagnetic field.

\section*{Part A. Free electron gas (1.4 points)}

Let's consider an infinite space filled with a metal: positively charged ions form a regular crystal lattice, free electrons move in the space between them, the concentrations of charges of both signs are the same and equal to $n$.

A uniform alternating electric field $\vec{E}\sin\omega t$ is created inside the metal. The ions are infinitely heavy and immobile, and the effective masses of the electrons are $m$. Within the framework of the simplest model, we can assume that the field acting on an individual electron coincides with $\vec{E}\sin\omega t$. All other forces (including dissipative ones) are small.

\problembox{A1}{The electric field leads to a consistent displacement $\vec{r}(t)$ of electrons in the direction of the field. Determine the dependence $\vec r(t)$ and the polarization vector of the medium $\vec P(t)$ in steady state.}{1.10}

\problembox{A2}{Calculate the dielectric constant of the metal $\varepsilon(\omega)$ and plot this dependence.}{0.30}

\section*{Part B. Plasmon resonance (7.9 points)}

Let us consider a dielectric ball with permittivity $\varepsilon$ in a uniform stationary electric field $\vec{E}$. In a uniform field, the moving charges inside the ball are displaced, and the field in the system can be represented as a superposition of the external field $\vec{E}$ and the field from two oppositely charged balls separated by a small distance along the vector $\vec{E}$ (see Fig. 1).

\begin{figure}[H]
    \centering
    \includegraphics[width=0.72\textwidth]{T1_EN_assets/fig1.jpg}
    \caption{Figure 1.}
\end{figure}

\problembox{B1}{Determine the field $\vec{E}_{in}$ inside the ball and its dipole moment $\vec{d}_0$.}{2.60}

Let us study the behavior of a metal ball of radius $R$ in the field of a plane wave with frequency $\omega$ and amplitude $\vec{E}_0$. It turns out that if the wavelength and depth of penetration of the field into the metal are much greater than the size of the ball, then its behavior is similar to the behavior of a dielectric ball in a constant electric field, only the value $\varepsilon(\omega)$, similar to that calculated in the previous paragraph, must be taken as the dielectric constant. As the wave passes, the amplitude of the external field at the location of the ball changes $\vec{E}=\vec{E}_0\cos\omega t$, and the dipole moment also changes $\vec{d}=\vec{d}_0\cos\omega t$.

\problembox{B2}{Draw a picture of the electric field lines in the system (inside, near and far from the ball), counting $\varepsilon(\omega)=-3$.}{0.70}

\problembox{B3}{At a certain value $\varepsilon(\omega_{res})=\varepsilon_{res}$ the amplitude of the electric field inside the ball $\left|\vec{E}_{in}\right|$ begins to increase without limit. Determine the value of $\varepsilon_{res}$.}{0.20}

A significant increase in the amplitude of the electric field inside and outside the ball at frequency $\omega_{res}$ is called plasmon resonance. We got an infinite answer because we did not take energy dissipation into account. Now let's take dissipation into account and assume that the main losses in the system are associated with dipole radiation.

\problembox{B4}{An oscillating dipole emits energy. Estimate the power $I$ of these losses using dimensional analysis and the fact that the intensity of the dipole radiation depends only on the amplitude of the dipole moment $\left|\vec{d}_0\right|$, its oscillation frequency $\omega_{res}$, the wave speed c and the electrical constant $\varepsilon_0$.}{0.60}

\problembox{B5}{Estimate the average energy pumped into the system due to the external field under plasmon resonance conditions at $\omega=\omega_{res}$. From the energy balance condition, estimate the amplitude of the field inside the ball in resonance.}{3.80}

\section*{Part C. Raman spectroscopy (1.3 points)}

Let's move on to the second part of the problem and consider Raman scattering, which also underlies SERS. Let's start the conversation by discussing the structure of molecules. We will imagine the molecule as a set of weights (atoms) connected to each other by springs (chemical bonds). For simplicity, in what follows we consider a molecule consisting of only two atoms.

\problembox{C1}{Let us consider two weights with masses $m_1$ and $m_2$, connected to each other by a spring with stiffness $k$. Determine the frequency $\omega_0$ of small oscillations in the system.}{0.70}

A complex molecule will be characterized not by one frequency $\omega_0$, but by a whole spectrum of natural frequencies. Knowledge of this spectrum allows us to unambiguously identify the substance with which we are dealing. This is the main idea of SERS. Let us study the behavior of a molecule in an external electric field $\vec{E}_0\cos\omega t$. For simplicity, we will assume that the atoms are not charged, i.e. in the absence of an external field, the molecule does not have a dipole moment. However, in an external field, like a metal ball, the molecule is polarized, i.e. acquires a dipole moment $$\vec{d} = \varepsilon_0\alpha\vec{E}_0\cos\omega t,$$ where $\alpha$ is the polarizability of the molecule. Here we assume that the direction of the induced dipole moment $\vec{d}$ coincides with the direction of the external field.

The mechanical vibrations studied in the previous paragraph never stop; the atoms in the molecule oscillate all the time. This means that the deviation of the distance between atoms from the equilibrium position $x$ changes with time, $x = x_0\cos\omega_0 t$. In turn, the ability of a molecule to polarize depends on the relative arrangement of atoms, that is, $\alpha = \alpha(x)$ (see Fig.~2).

\begin{figure}[H]
    \centering
    \includegraphics[width=0.72\textwidth]{T1_EN_assets/fig2.jpg}
    \caption{Figure 2.}
\end{figure}

\problembox{C2}{Considering the amplitude of mechanical vibrations near the equilibrium position to be small, write an explicit expression for $\alpha(x)$ in the linear approximation, if $\alpha(0)=\alpha_0$ and $\frac{d\alpha}{dx}\bigg|_{x=0}=\beta_0$.}{0.20}

\problembox{C3}{Calculate the dipole moment of a diatomic molecule in an external field $\vec{E}_0\cos\omega t$. Present the answer as: $$\vec{d}=\sum_i\vec{d}_i\cos\omega_it $$}{0.20}

\problembox{C4}{The detector registers the dipole radiation of the molecule. The expression from the previous paragraph will correspond to several peaks at frequencies $\omega_i$. The height of each peak is equal to the emission intensity of an individual dipole $\vec{d}_i$. For each peak, determine its frequency and height. Express your answer in terms of $\varepsilon_0$, $\alpha_0$, $\beta_0$, $x_0$, $\omega$, $\omega_0$ and $\vec{E}_0$.}{0.20}

The presence of peaks at frequencies other than $\omega$ is called Raman scattering. In order for the detector to receive a good signal, it is necessary to place the molecule in a strong electric field. This strong field can be created using plasmon resonance. This is the main difference between the idea of SERS and conventional Raman spectroscopy.

\section*{Part D. Surface-enhanced Raman spectroscopy (3.4 points)}

Let us consider a ball with dielectric constant $\varepsilon(\omega)$ in a uniform alternating field of amplitude $\left|\vec{E}_0\right|$ under plasmon resonance conditions.

\problembox{D1}{At what point in space should the molecule of the substance under study be placed so that it is in the electric field of maximum strength?}{1.20}

\problembox{D2}{Calculate the electric field enhancement factor $g(\omega)$ at this point, $g(\omega)=\underset{\vec{r}}{\max}\frac{\left|\vec{E}(\vec{r})\right|}{\left|\vec{E}_0\right|}$. Express your answer in terms of the dielectric constant of the metal $\varepsilon(\omega)$.}{0.30}

The metal particle amplifies not only the original radiation with amplitude $\left|\vec{E}_0\right|$, but also the dipole radiation that the molecule generates. This process can also be characterized by the amplification factor $g'(\omega, \omega_0)$. If the frequency $\omega$ of the external field is large compared to the frequency $\omega_0$ of the natural vibrations of the molecule, we can assume $g'\approx{g}$. Then the signal intensity in the SERS experiment increases compared to conventional Raman spectroscopy by $g^4$ times.

In real experiments, to amplify a signal, they usually study not just one molecule, but a large number of them at once. The radiation from different molecules in such experiments is not coherent, therefore the total intensity of radiation from $N$ molecules is equal to $NI_0$, where $I_0$ is the intensity of radiation from one molecule. An example of the results of such measurements is shown in Fig. 3.

\begin{figure}[H]
    \centering
    \includegraphics[width=0.72\textwidth]{T1_EN_assets/fig3.jpg}
    \caption{Figure 3.}
\end{figure}

Spectrum of the Raman signal. The upper curve corresponds to the SERS experiment, the lower curve corresponds to conventional Raman spectroscopy. The horizontal axis represents the magnitude of the Raman shift, $k_0=\omega_0/c$, and the vertical axis represents the radiation intensity in arbitrary units. Please note that in these two experiments a different number of molecules participated in the formation of the signal, see the captions on the graph.

\problembox{D3}{From the experimental data in the figure above, estimate the electric field enhancement factor $g$ due to plasmon resonance (consider the peak at the Raman shift $\omega_0/c=1000~\text{cm}^{-1}$). Consider that condition $\omega_0\ll{\omega}$ is satisfied.}{1.40}

\problembox{D4}{Using the results of the second part of the problem, estimate the size $R$ of the metal particles used in the experiment if the wavelength of the incident radiation is $\lambda=785~\text{nm}$.}{0.50}

\vspace{1em}
\noindent\textit{Source:} \url{https://pho.rs/p/3042}

\end{document}